% problem_id: S001-proposition-finite-type-injective-into-flat-mod-m
% source_id: S001 (Stacks Project)
% source_file: flat.tex
% statement_locator: flat.tex lines 7413--7421
% proof_locator: flat.tex lines 7423--7549
% env: proposition
% status: text-extracted-verbatim (difficulty judgement pending, written by hand)
%
% The statement and proof below are the author's own text, extracted verbatim.
% Nothing is paraphrased, shortened, or supplemented.

\begin{proposition}
\label{proposition-finite-type-injective-into-flat-mod-m}
Let $A \to B$ be a local ring homomorphism of local rings
which is essentially of finite type. Let $M$ be a flat $A$-module,
$N$ a finite $B$-module and $u : N \to M$ an $A$-module map such that
$\overline{u} : N/\mathfrak m_AN \to M/\mathfrak m_AM$ is injective.
Then $u$ is $A$-universally injective, $N$ is of finite presentation over
$B$, and $N$ is flat over $A$.
\end{proposition}

\begin{proof}
We may assume that $B$ is the localization of a finitely presented
$A$-algebra $B_0$ and that $N$ is the localization of a
finitely presented $B_0$-module $M_0$, see
Lemma \ref{lemma-reduce-finite-type-injective-into-flat-mod-m}.
By More on Morphisms, Lemma
\ref{more-morphisms-lemma-generic-flatness-stratification}
there exists a ``generic flatness stratification''
for $\widetilde{M_0}$ on $\Spec(B_0)$ over $\Spec(A)$.
Translating back to $N$ we find a sequence of closed subschemes
$$
S = \Spec(A) \supset S_0 \supset S_1 \supset \ldots \supset S_t = \emptyset
$$
with $S_i \subset S$ cut out by a finitely generated ideal of $A$
such that the pullback of $\widetilde{N}$ to
$\Spec(B) \times_S (S_i \setminus S_{i + 1})$ is flat over
$S_i \setminus S_{i + 1}$. We will prove the proposition by
induction on $t$ (the base case $t = 1$ will be proved in parallel
with the other steps). Let $\Spec(A/J_i)$ be the scheme theoretic
closure of $S_i \setminus S_{i + 1}$.

\medskip\noindent
{\bf Claim 1.} $N/J_iN$ is flat over $A/J_i$. This is immediate for
$i = t - 1$ and follows from the induction hypothesis for $i > 0$.
Thus we may assume $t > 1$, $S_{t - 1} \not = \emptyset$, and
$J_0 = 0$ and we have to prove that $N$ is flat. Let $J \subset A$
be the ideal defining $S_1$. By induction on $t$ again, we also
have flatness modulo powers of $J$. Let $A^h$ be the henselization of $A$
and let $B'$ be the localization of $B \otimes_A A^h$ at the maximal ideal
$\mathfrak m_B \otimes A^h + B \otimes \mathfrak m_{A^h}$. Then $B \to B'$
is faithfully flat. Set $N' = N \otimes_B B'$. Note that $N'$
is $A^h$-flat if and only if $N$ is $A$-flat. By
Theorem \ref{theorem-flattening-local} there is a smallest ideal
$I \subset A^h$ such that $N'/IN'$ is flat over $A^h/I$, and
$I$ is finitely generated. By the above $I \subset J^nA^h$ for
all $n \geq 1$. Let $S_i^h \subset \Spec(A^h)$ be the inverse image
of $S_i \subset \Spec(A)$. By Lemma \ref{lemma-closed-points-complement}
we see that $V(I)$ contains the closed points of $U = \Spec(A^h) - S_1^h$.
By construction $N'$ is $A^h$-flat over $U$.
By Lemma \ref{lemma-make-smaller-flatness-ideal} we see that $N'/I_2N'$
is flat over $A/I_2$, where
$I_2 = \Ker(I \to \Gamma(U, I/I^2))$. Hence $I = I_2$ by minimality
of $I$. This implies that $I = I^2$ locally on $U$, i.e.,
we have $I\mathcal{O}_{U, u} = (0)$ or $I\mathcal{O}_{U, u} = (1)$
for all $u \in U$. Since $V(I)$ contains the closed points of $U$
we see that $I = 0$ on $U$. Since $U \subset \Spec(A^h)$ is scheme
theoretically dense (because replaced $A$ by $A/J_0$ in the beginning
of this paragraph), we see that $I = 0$. Thus $N'$ is $A^h$-flat
and hence Claim 1 holds.

\medskip\noindent
We return to the situation as laid out before Claim 1. With
$A^h$ the henselization of $A$, with $B'$ the localization
of $B \otimes_A A^h$ at the maximal ideal
$\mathfrak m_B \otimes A^h + B \otimes \mathfrak m_{A^h}$, and with
$N' = N \otimes_B B'$ we now see that the flattening ideal $I \subset A^h$
of Theorem \ref{theorem-flattening-local} is nilpotent.
If $nil(A^h)$ denotes the ideal of nilpotent elements, then
$nil(A^h) = nil(A) A^h$
(More on Algebra, Lemma \ref{more-algebra-lemma-henselization-nil}).
Hence there exists a finitely generated
nilpotent ideal $I_0 \subset A$ such that $N/I_0N$ is flat over $A/I_0$.

\medskip\noindent
{\bf Claim 2.} For every prime ideal $\mathfrak p \subset A$
the map
$\kappa(\mathfrak p) \otimes_A N \to \kappa(\mathfrak p) \otimes_A M$
is injective. We say $\mathfrak p$ is bad it this is false.
Suppose that $C$ is a nonempty chain of bad primes and set
$\mathfrak p^* = \bigcup_{\mathfrak p \in C} \mathfrak p$.
By Lemma \ref{lemma-find-pure-spreadout}
there is a finitely generated ideal
$\mathfrak a \subset \mathfrak p^*A_{\mathfrak p^*}$
such that there is a pure spreadout over $V(\mathfrak a)$.
If $\mathfrak p^*$ were good, then it would follow from
Lemma \ref{lemma-properties-pure-spreadout}
that the points of $V(\mathfrak a)$ are good.
However, since $\mathfrak a$ is finitely generated and since
$\mathfrak p^*A_{\mathfrak p^*} = \bigcup_{\mathfrak p \in C}A_{\mathfrak p^*}$
we see that $V(\mathfrak a)$ contains a $\mathfrak p \in C$, contradiction.
Hence $\mathfrak p^*$ is bad. By Zorn's lemma, if there exists a
bad prime, there exists a maximal one, say $\mathfrak p$.
In other words, we may assume every $\mathfrak p' \supset \mathfrak p$,
$\mathfrak p' \not = \mathfrak p$ is good.
In this case we see that for every $f \in A$, $f \not \in \mathfrak p$
the map $u \otimes \text{id}_{A/(\mathfrak p + f)}$ is universally
injective, see Lemma \ref{lemma-universally-injective-if-flat}.
Thus it suffices to show that $N/\mathfrak p N$ is separated
for the topology defined by the submodules $f(N/\mathfrak pN)$.
Since $B \to B'$ is faithfully flat, it is enough to prove the
same for the module $N'/\mathfrak p N'$.
By Lemma \ref{lemma-flat-finite-type-local-colimit-free} and
More on Algebra, Lemma
\ref{more-algebra-lemma-content-exists-flat-Mittag-Leffler}
elements of $N'/\mathfrak pN'$ have content ideals in $A^h/\mathfrak pA^h$.
Thus it suffices to show that
$\bigcap_{f \in A, f \not \in \mathfrak p} f(A^h/\mathfrak p A^h) = 0$.
Then it suffices to show the same for
$A^h/\mathfrak q A^h$ for every prime $\mathfrak q \subset A^h$
minimal over $\mathfrak p A^h$. Because $A \to A^h$ is the henselization,
every $\mathfrak q$ contracts to $\mathfrak p$ and every
$\mathfrak q' \supset \mathfrak q$, $\mathfrak q' \not = \mathfrak q$
contracts to a prime $\mathfrak p'$ which strictly contains $\mathfrak p$.
Thus we get the vanishing of the intersections from
Lemma \ref{lemma-big-intersection-is-zero}.

\medskip\noindent
At this point we can put everything together. Namely, using
Claim 1 and Claim 2 we see that $N/I_0 N \to M/I_0M$ is
$A/I_0$-universally injective by
Lemma \ref{lemma-universally-injective-if-flat}.
Then the diagrams
$$
\xymatrix{
N \otimes_A (I_0^n/I_0^{n + 1}) \ar[r] \ar[d] &
M \otimes_A (I_0^n/I_0^{n + 1}) \ar@{=}[d] \\
I_0^n N /I_0^{n + 1} N \ar[r] &
I_0^n M /I_0^{n + 1} M
}
$$
show that the left vertical arrows are injective. Hence by
Algebra, Lemma \ref{algebra-lemma-what-does-it-mean-again}
we see that $N$ is flat. In a similar way the
universal injectivity of $u$ can be reduced (even without
proving flatness of $N$ first) to the one modulo $I_0$. This finishes
the proof.
\end{proof}
